\section{Ultraviolet--infrared hierarchy}

\subsection{UVIR-001: the minimal condensate route closes negative}

The tested microscopic candidate is
\begin{equation}
 S_\Phi=\int \dd^4x\sqrt{-g}\left[
 -\nabla_\mu\Phi^*\nabla^\mu\Phi
 -m^2|\Phi|^2
 -\frac{\lambda_4}{2}|\Phi|^4
 -\frac{\lambda_6}{3\Lambda^2}|\Phi|^6\right].
 \label{eq:candidate-uv}
\end{equation}
Writing $\Phi=\rho e^{i\Theta}/\sqrt2$ and
$Z=-\nabla_\mu\Theta\nabla^\mu\Theta$, its stable finite-density branch can
be reduced to a phase action $P(Z)$.  In the pure-sextic limit the explicit
amplitude integration gives
\begin{equation}
 P(Z)=\frac{2\Lambda}{3\sqrt{\lambda_6}}(Z-m^2)^{3/2}.
 \label{eq:sextic-pz}
\end{equation}
This nonanalytic power does not supply the required static spatial operator.
About a timelike background $Z_0=\mu^2$, a static phase gradient
$q=|\nabla\vartheta|$ instead yields
\begin{equation}
 P(\mu^2-q^2)-P(\mu^2)=-P_Z(\mu^2)q^2+\mathcal O(q^4),
 \qquad P_Z=\frac{\rho_0^2}{2}>0.
 \label{eq:uvir001-static}
\end{equation}
The leading response is therefore quadratic for every stable nonzero branch
with the canonical phase kinetic term.  UVIR-001 is closed negative for this
minimal candidate.  The result rejects neither the condensate's density and
topology roles nor a bottom-up spatial $Y^{3/2}$ EFT.

\subsection{UVIR-002: provisional two-sector route}

UVIR-002 also rejects a standalone nonrelativistic gradient-dominated branch:
the branch producing the desired cubic spatial action has a negative quadratic
time kinetic coefficient.  The least-incomplete route tested separates the
jobs of the condensate and force mediator.  The independent invariants are
\begin{equation}
 \mathcal Q=U^\mu\nabla_\mu\psi,
 \qquad
 \mathcal Y=h^{\mu\nu}\nabla_\mu\psi\nabla_\nu\psi.
 \label{eq:preferred-invariants}
\end{equation}
Here $\Phi$ retains finite density, circulation, defects and topology, while
$(\psi,U^\mu)$ carries the conditional low-acceleration response.  This is a
provisional architecture selection, not microscopic closure.

\subsection{UVIR-003 Stage A action and limited pass}

Stage A chooses an independent unit vector and retains the complete
Einstein--aether two-derivative basis,
\begin{align}
 \mathcal L_U=-\frac{M_U^2}{2}\big[&c_1(\nabla_\mu U_\nu)(\nabla^\mu U^\nu)
 +c_2(\nabla_\mu U^\mu)^2 \nonumber\\
 &+c_3(\nabla_\mu U_\nu)(\nabla^\nu U^\mu)-c_4a_\mu a^\mu\big]
 +\frac{\lambda}{2}(U^\mu U_\mu+1),
 \label{eq:aether-action}
\end{align}
where $a_\mu=U^\nu\nabla_\nu U_\mu$.  A positive alignment energy couples
the frame to the condensate current without imposing equality.  The regulated
force truncation is
\begin{equation}
 \mathcal L_\psi=\frac{K_Q}{2}\mathcal Q^2-A\mathcal Y^{3/2}
 -\frac{\gamma}{2M_*^2}(\Delta_U\psi)^2.
 \label{eq:stage-a-force}
\end{equation}
The symbolic Stage-A calculation gives the necessary frozen-metric conditions
\begin{equation}
 c_{14}>0,\qquad c_1>0,\qquad c_{123}>0,
 \qquad K_Q>0,\quad A>0,\quad\gamma>0,
 \label{eq:stage-a-signs}
\end{equation}
and at zero background gradient the regulator produces
\begin{equation}
 \omega^2=\frac{\gamma}{K_QM_*^2}k^4.
 \label{eq:k4-dispersion}
\end{equation}
This is a pass only in the declared flat, constant-frame decoupling limit.

\subsection{Stage B frame-sector diagnostic}

A verified sign map to the Einstein--aether review
\cite{elingjacobsonmattingly2004} preserves all four operator labels.  Its
normalization also requires
\begin{equation}
 \alpha_i=r_Uc_i,\qquad r_U=\frac{M_U^2}{M_P^2},
 \label{eq:aether-normalization}
\end{equation}
because the source places the aether operators under the Einstein--Hilbert
prefactor, while the Stage-A action declares an independent $M_U$.  The
published coupled metric--aether result therefore gives
\begin{align}
 s_2^2&=\frac{1}{1-\alpha_{13}},\\
 s_1^2&=\frac{\alpha_1-\tfrac12\alpha_1^2+\tfrac12\alpha_3^2}
 {\alpha_{14}(1-\alpha_{13})},\\
 s_0^2&=\frac{\alpha_{123}(2-\alpha_{14})}
 {\alpha_{14}(1-\alpha_{13})(2+\alpha_{13}+3\alpha_2)}.
 \label{eq:aether-speeds}
\end{align}
As $r_U\to0$ at fixed $c_i$ ratios, the spin-1 and spin-0 expressions
reduce to the Stage-A ratios $c_1/c_{14}$ and $c_{123}/c_{14}$.  This confirms
that those ratios are weak-metric-coupling limits, not the full mode speeds.
It is a checked literature substitution, not an independent ITSM-side
linearized Einstein-equation derivation, and it excludes the coupled
$\Phi$--$U$--$\psi$ system.

\subsection{Zero-gradient force-block factorization}

On the declared background $\bar\psi=\text{constant}$, every derivative of
the force field begins at first perturbative order.  Consequently,
\begin{equation}
 \mathcal Q^2=\dot\pi^2+\mathcal O(3),\qquad
 \mathcal Y^{3/2}=\mathcal O(3),\qquad
 (\Delta_U\psi)^2=(\nabla^2\pi)^2+\mathcal O(3),
\end{equation}
where metric and frame corrections enter only in the indicated higher-order
terms.  The force scalar therefore factorizes from the lapse, shift, frame and
condensate blocks at quadratic order for the declared truncation:
\begin{equation}
 S_\pi^{(2)}=\frac12\int\dd^4x\left[
 K_Q\dot\pi^2-\frac{\gamma}{M_*^2}(\nabla^2\pi)^2\right].
 \label{eq:zero-gradient-force-block}
\end{equation}
It carries one nonnegative $z=2$ scalar mode for $K_Q,\gamma>0$.  This bounded
result does not reduce the remaining metric--frame--condensate block.

A constant field redefinition $\psi_c=s\psi$ sends
\begin{equation}
 K_Q\to\frac{K_Q}{s^2},\quad
 A\to\frac{A}{s^3},\quad
 \gamma\to\frac{\gamma}{s^2},\quad
 C_m\to\frac{C_m}{s},\quad q\to sq.
\end{equation}
Thus $\gamma/K_Q$, $A/K_Q^{3/2}$, $C_m/\sqrt{K_Q}$ and $Aq/K_Q$ are invariant,
whereas $K_Q$ alone is not.  UVIR-003 can determine a structural stability
domain in invariant ratios, but a parent microscopic calculation or MAT-001
must fix the physical normalization and select a point in that domain.  The
earlier $K_Q\sim M_P^2$ estimate remains a speculative dimensional analogy.
\subsection{Scalar ADM readiness}

The intended coupled reduction was to be performed around Minkowski space
with $\Phi=\rho_0e^{i\mu t}/\sqrt2$, constant $U^\mu$ and constant
$\bar\psi$.  This background is not on shell for the declared gravitational
action.  Writing $s=\rho_0^2$, the homogeneous condensate has
\begin{equation}
 p_\Phi=\frac12s\mu^2-V(\rho_0),\qquad
 \rho_\Phi=\frac12s\mu^2+V(\rho_0),\qquad
 \rho_\Phi+p_\Phi=s\mu^2>0.
 \label{eq:condensate-enthalpy}
\end{equation}
A constant vacuum-energy subtraction shifts $\rho$ and $p$ oppositely, so it
cannot cancel their sum.  It therefore cannot make both flat-background
Einstein equations vanish on the nonzero finite-density branch.

An exact stationary completion would need a reservoir or driver satisfying
\begin{equation}
 \rho_R=-\rho_\Phi,\qquad p_R=-p_\Phi,\qquad
 \rho_R+p_R=-s\mu^2.
 \label{eq:required-support-stress}
\end{equation}
This is non-vacuum stress.  Its scalar perturbations generally enter the
lapse and shift constraints and cannot be silently omitted.  Until a
covariant support action, a controlled rigid-support limit, or a consistent
cosmological background is declared, eliminating the ADM constraints would
produce an off-shell kinetic matrix.  The scalar reduction is therefore
blocked at readiness rather than counted as a stability failure.  The
zero-gradient force factorization remains valid because that block is
quadratically independent on the declared constant-force background.

\subsection{Background-completion screen}

For a minimally coupled shift-symmetric support scalar with
$X_R=-\nabla_\mu\chi\nabla^\mu\chi/2>0$ and
$\mathcal L_R=P_R(X_R)$,
\begin{equation}
 \rho_R+p_R=2X_RP_{R,X}.
 \label{eq:px-support-enthalpy}
\end{equation}
Exact Minkowski support would therefore require
\begin{equation}
 P_{R,X}=-\frac{s\mu^2}{2X_R}<0.
 \label{eq:px-support-sign}
\end{equation}
The principal quadratic fluctuation action is
\begin{equation}
 \mathcal L_R^{(2)}
 =\frac12(P_{R,X}+2X_RP_{R,XX})\dot\pi_R^2
 -\frac12P_{R,X}(\nabla\pi_R)^2.
 \label{eq:px-support-quadratic}
\end{equation}
A healthy two-derivative short-wavelength scalar requires both
$P_{R,X}+2X_RP_{R,XX}>0$ and $P_{R,X}>0$, in conflict with
Eq.~\eqref{eq:px-support-sign}.  This is the local version of the familiar
instability obstruction for simple phantom-like scalar backgrounds
\cite{vikman2005}.  The ghost-condensate point has $P_{R,X}=0$ and hence zero
enthalpy; its $k^4$ excitation does not by itself supply the missing
background stress \cite{arkani-hamed2004}.

Stable NEC violation can be engineered with additional EFT structure
\cite{creminelli2006}, but importing that structure would add a new
microscopic sector, cutoff and stability gate.  A prescribed rigid stress is
also insufficient for the full problem because it has no action-derived
lapse or shift response.  The screen therefore selects no new exact-Minkowski
support field.

The least-assumptive route is a self-consistent evolving flat-FRW background,
\begin{equation}
 M_{\rm cos}^2=M_P^2+\frac{M_U^2}{2}(c_1+3c_2+c_3),
 \label{eq:cosmological-planck-mass}
\end{equation}
\begin{equation}
 3M_{\rm cos}^2H^2=\rho_{\rm total},\qquad
 -2M_{\rm cos}^2\dot H=\rho_{\rm total}+p_{\rm total}.
 \label{eq:selected-frw-background}
\end{equation}
For isolated condensate charge $n=\mu\rho^2$,
\begin{equation}
 \dot n+3Hn=0.
 \label{eq:condensate-charge-dilution}
\end{equation}
Thus the amplitude or chemical potential generally evolves.  Exact constant
$n$ in expansion would require a separately declared charge source
$S_N=3Hn$; the stress-energy exchange vector $Q_{\rm syn}^{\nu}$ is not
automatically that source.

\subsection{Evolving FRW existence check}

For the comoving unit aether, the Stage-A contractions reduce to
$3H_N^2$, $9H_N^2$, $3H_N^2$ and zero acceleration squared.  The alignment
projector annihilates the parallel homogeneous condensate current, and the
constant force field has no background derivatives.  The reduced action is
\begin{equation}
 S_{\rm bg}=\int\dd t\left[
 -\frac{3M_{\rm cos}^2a\dot a^2}{N}
 +\frac{a^3}{2N}(\dot\rho^2+\rho^2\dot\Theta^2)
 -Na^3V(\rho)\right].
 \label{eq:frw-minisuperspace-action}
\end{equation}
Its amplitude and phase equations are
\begin{equation}
 \ddot\rho+3H\dot\rho-\rho\mu^2+V_{,\rho}=0,
 \qquad
 \frac{\dd}{\dd t}(a^3\rho^2\mu)=0.
 \label{eq:frw-condensate-equations}
\end{equation}

A representative dimensionless integration, not a parameter fit, uses
$M_P=1$, $M_U=0.5$, $(c_1,c_2,c_3,c_4)=(0.10,0.05,0.05,0.05)$ and positive
mass--quartic--sextic couplings.  From $t=0$ to $8$, it remains regular and
expanding:
\begin{equation}
 (a,\rho,\mu,H):
 (1,1,1.123610,0.619841)
 \longrightarrow
 (4.571577,0.134236,0.652648,0.076405).
 \label{eq:frw-existence-endpoints}
\end{equation}
The maximum relative Friedmann residual is $2.124\times10^{-10}$, charge drift
is $2.220\times10^{-16}$ and the continuity residual is
$1.898\times10^{-15}$.  This verifies an on-shell existence branch and removes
the background blocker.  It does not select a fiducial cosmology or establish
perturbation stability.

The metric--frame--condensate constraint reduction is now ready to begin on
this time-dependent background.  A subhorizon limit can test local modes only for
$k/a\gg H$; the strict low-$k$ question remains part of the complete
cosmological perturbation system.

\subsection{Scalar ADM principal-symbol reduction}

The first scalar constraint subgate uses aether-unitary gauge,
\begin{equation}
 N=1+\delta N,\qquad N_i=\partial_i\beta,\qquad
 h_{ij}=a^2e^{2\mathcal R}\delta_{ij},\qquad U^\mu=n^\mu,
 \label{eq:scalar-adm-gauge}
\end{equation}
and freezes the evolving background coefficients over a wavelength.  For
$q_{\rm phys}=k/a\gg H$, the principal lapse and scalar-shift solutions are
\begin{equation}
 \delta N=-\frac{2M_P^2}{M_U^2c_{14}}\mathcal R
 +\frac{\dot\rho\,\delta\dot\rho+\rho^2\mu\dot\vartheta}
 {M_U^2c_{14}q_{\rm phys}^2},
 \label{eq:scalar-adm-lapse}
\end{equation}
\begin{equation}
 \beta=-\frac{2M_{\rm cos}^2\dot{\mathcal R}
 +\dot\rho\,\delta\rho+\rho^2\mu\vartheta}
 {M_U^2c_{123}q_{\rm phys}^2}.
 \label{eq:scalar-adm-shift}
\end{equation}
With $F=M_{\rm cos}^2/M_P^2
=1+(\alpha_{13}+3\alpha_2)/2$, their elimination gives
\begin{equation}
 \mathcal L_{\mathcal R}^{(2)}
 =\mathcal K_{\mathcal R}\dot{\mathcal R}^2
 -\mathcal G_{\mathcal R}q_{\rm phys}^2\mathcal R^2,
 \quad
 \mathcal K_{\mathcal R}
 =\frac{2M_P^2F(1-\alpha_{13})}{\alpha_{123}},
 \quad
 \mathcal G_{\mathcal R}
 =\frac{M_P^2(2-\alpha_{14})}{\alpha_{14}}.
 \label{eq:scalar-adm-reduced}
\end{equation}
Their ratio independently recovers the exact coupled Einstein--aether spin-0
speed \cite{elingjacobsonmattingly2004},
\begin{equation}
 s_0^2=
 \frac{\alpha_{123}(2-\alpha_{14})}
 {\alpha_{14}(1-\alpha_{13})(2+\alpha_{13}+3\alpha_2)}.
 \label{eq:scalar-adm-spin-zero}
\end{equation}

The condensate velocity Hessian retains a leading lapse-induced correction
whose determinant factorizes as
\begin{equation}
 \det\mathbf K_\Phi(q_{\rm phys})
 =\rho^2\left[
 1-\frac{\dot\rho^2+\rho^2\mu^2}
 {M_U^2c_{14}q_{\rm phys}^2}\right].
 \label{eq:scalar-adm-condensate-determinant}
\end{equation}
Thus the controlled domain is not merely $q_{\rm phys}\gg H$ but
\begin{equation}
 q_{\rm phys}\gg\max(H,q_{\rm ADM}),\qquad
 q_{\rm ADM}^2=
 \frac{\dot\rho^2+\rho^2\mu^2}{M_U^2c_{14}}.
 \label{eq:scalar-adm-validity}
\end{equation}
At the representative dimensionless point,
$\mathcal K_{\mathcal R}=39.94375$,
$\mathcal G_{\mathcal R}=52.33333$ and $s_0^2=1.31018$.  The principal block
is positive, but its cone is wider than the metric cone.  This is a bounded
subhorizon result with a multicone-causality flag, not a full finite-$k$ or
low-$k$ cosmological stability proof.

\subsection{Causality and cutoff status}

Around a nonzero spatial background gradient of magnitude $q$, the force
fluctuation has
\begin{equation}
 \omega^2=\frac{3Aq(1+\cos^2\theta)}{K_Q}k^2
 +\frac{\gamma}{K_QM_*^2}k^4.
 \label{eq:force-dispersion}
\end{equation}
Thus the long-wavelength metric-light-cone threshold depends on the unmatched
ratio $Aq/K_Q$, while phase and group velocities grow without bound at large
$k$ inside the truncated dispersion.  Preferred-frame metric superluminality
is not by itself a proof of acausality, but the foliation-causality test and the
physical EFT domain have not been established.

Expansion of the cubic vertex supplies only the restricted, time-normalized
longitudinal IR NDA scale
\begin{equation}
 \Lambda_{\rm NDA,long}^{\rm (IR)}\sim\frac{K_Q^{3/4}}{\sqrt A}.
 \label{eq:long-nda}
\end{equation}
It is not the physical cutoff $\Lambda_{\rm EFT}$: transverse normalization,
$k^4$-dominated power counting and a loop or unitarity analysis remain open.
The comparison of the light-cone crossing scale with $\Lambda_{\rm EFT}$
therefore cannot yet be made.

A separate dimensional analogy, $K_Q=k_QM_P^2$, would place the
long-wavelength threshold at order $a_0$ for order-unity coefficients.  That
exercise is explicitly speculative: it is introduced only as a priority
signal for deriving $K_Q$, not as a property of ITSM.

\subsection{Conditional infrared operator}

Define the normalized spatial invariant
\begin{equation}
 Y=\frac{h^{\mu\nu}\nabla_\mu\psi\nabla_\nu\psi}{a_0^2}.
 \label{eq:Y}
\end{equation}
The target low-acceleration operator remains
\begin{equation}
 \mathcal L_{\psi,\rm IR}
 =-\frac{2C_{\rm IR}}{3}\mpl^2a_0^2Y^{3/2}
 +\mathcal O(Y^2,\nabla^2/\Lambda_{\rm EFT}^2).
 \label{eq:ir-action}
\end{equation}
Neither $C_{\rm IR}$, $K_Q$ nor the normalized matter vertex has been derived.
The preceding analytic Born--Infeld expression also remains rejected as the
source of this branch because its small-invariant equation is linear at
leading order.  UVIR-003 is therefore in progress and MAT-001 remains blocked.
